\documentclass[11pt,letterpaper]{article}

% ── Geometry ──────────────────────────────────────────────────────────────────
\usepackage[margin=1in]{geometry}

% ── Pagination ────────────────────────────────────────────────────────────────
\usepackage[all]{nowidow}    % Require ≥2 lines on each side of a page break
\brokenpenalty=10000         % No hyphenated words across page breaks
\raggedbottom                % Allow uneven page bottoms rather than stretching

% ── Fonts ─────────────────────────────────────────────────────────────────────
\usepackage[T1]{fontenc}
\usepackage{newpxtext}     % Palatino-based serif (TeX Gyre Pagella)
\usepackage{newpxmath}     % Matching math font

% ── Colors ────────────────────────────────────────────────────────────────────
\usepackage[dvipsnames]{xcolor}
\definecolor{SectionBlue}{HTML}{1B3A5C}
\definecolor{AccentGray}{HTML}{4A4A4A}
\definecolor{QuoteBorder}{HTML}{8B9DAF}
\definecolor{RiskRed}{HTML}{C0392B}
\definecolor{RiskAmber}{HTML}{E67E22}
\definecolor{RiskGreen}{HTML}{27AE60}
\definecolor{BoxBg}{HTML}{F7F8FA}

% ── Tables ────────────────────────────────────────────────────────────────────
\usepackage{booktabs}
\usepackage{tabularx}

% ── Boxes ─────────────────────────────────────────────────────────────────────
\usepackage[framemethod=tikz]{mdframed}

% ── Lists ─────────────────────────────────────────────────────────────────────
\usepackage{enumitem}
\usepackage{needspace}

% ── Hyperlinks ────────────────────────────────────────────────────────────────
\usepackage{hyperref}
\hypersetup{
  colorlinks=true,
  linkcolor=SectionBlue,
  urlcolor=SectionBlue,
  citecolor=SectionBlue,
  pdfauthor={JM Freeman},
  pdftitle={vox --- PR/FAQ},
  pdfsubject={Working Backwards --- Product Discovery},
}

% ── Bibliography ─────────────────────────────────────────────────────────────
\usepackage[style=numeric,backend=biber,sorting=none]{biblatex}
\addbibresource{prfaq.bib}

% ── Section styling ───────────────────────────────────────────────────────────
\usepackage{titlesec}
\titleformat{\section}
  {\needspace{6\baselineskip}\Large\bfseries\color{SectionBlue}}
  {}
  {0pt}
  {}

\titleformat{\subsection}
  {\needspace{5\baselineskip}\large\bfseries\color{AccentGray}}
  {}
  {0pt}
  {}

\titleformat{\subsubsection}
  {\normalsize\bfseries\color{AccentGray}}
  {}
  {0pt}
  {}
\titlespacing*{\subsubsection}{0pt}{1em}{0.5em}[0pt]

% ── Document stage ───────────────────────────────────────────────────────────
\newcommand{\prfaqstage}[1]{\def\theprfaqstage{#1}}
\prfaqstage{shipped}

% ── Document version ─────────────────────────────────────────────────────────
\newcommand{\prfaqversion}[2]{\def\theprfaqmajor{#1}\def\theprfaqminor{#2}}
\prfaqversion{2}{2}

% ── Header/footer ────────────────────────────────────────────────────────────
\usepackage{fancyhdr}
\pagestyle{fancy}
\fancyhf{}
\fancyhead[L]{\footnotesize\textcolor{AccentGray}{Stage: \theprfaqstage{} \enspace\textbar\enspace v\theprfaqmajor.\theprfaqminor}}
\fancyhead[R]{\footnotesize\textcolor{AccentGray}{\thepage}}
\renewcommand{\headrulewidth}{0pt}

% ══════════════════════════════════════════════════════════════════════════════
% Custom environments
% ══════════════════════════════════════════════════════════════════════════════

% ── \prsection{title} — styled press release section header ──────────────────
\newcommand{\prsection}[1]{%
  \vspace{1em}%
  {\large\bfseries\color{SectionBlue}#1}%
  \vspace{0.3em}\par%
}

% ── customerquote — indented, italic customer quote ──────────────────────────
\newmdenv[
  topline=false,
  bottomline=false,
  rightline=false,
  linewidth=3pt,
  linecolor=QuoteBorder,
  backgroundcolor=BoxBg,
  innerleftmargin=12pt,
  innerrightmargin=12pt,
  innertopmargin=8pt,
  innerbottommargin=8pt,
  skipabove=\baselineskip,
  skipbelow=\baselineskip,
]{customerquote}

% ── spokespersonquote — indented spokesperson quote ──────────────────────────
\newmdenv[
  topline=false,
  bottomline=false,
  rightline=false,
  linewidth=3pt,
  linecolor=SectionBlue,
  backgroundcolor=BoxBg,
  innerleftmargin=12pt,
  innerrightmargin=12pt,
  innertopmargin=8pt,
  innerbottommargin=8pt,
  skipabove=\baselineskip,
  skipbelow=\baselineskip,
]{spokespersonquote}

% ── faqpair — numbered Q&A pair with bold question ─────────────────────────
\newcounter{faqnum}
\newenvironment{faqpair}[1]{%
  \needspace{4\baselineskip}%
  \vspace{0.5em}%
  \refstepcounter{faqnum}%
  \noindent\textbf{\color{SectionBlue}Q\thefaqnum: #1}\par\nopagebreak[4]%
  \vspace{0.2em}%
  \begin{adjustwidth}{1em}{0pt}%
  \setlength{\parindent}{0pt}%
}{%
  \end{adjustwidth}%
  \vspace{0.5em}%
}
% ── \faqref{faq:slug} — clickable "FAQ 3" reference ──────────────────────────
\newcommand{\faqref}[1]{\hyperref[#1]{FAQ~\ref*{#1}}}

% ── \featureitem — numbered, referenceable feature entry ──────────────────────
\newcounter{featurenum}
\newcommand{\featureitem}[2]{%
  \refstepcounter{featurenum}%
  \item[\textbf{\color{SectionBlue}F\thefeaturenum.}]%
  \textbf{#1} --- #2%
}
\newcommand{\featureref}[1]{\hyperref[#1]{Feature~\ref*{#1}}}

% ── adjustwidth (for faqpair indentation) ────────────────────────────────────
\usepackage{changepage}

% ══════════════════════════════════════════════════════════════════════════════
% Document
% ══════════════════════════════════════════════════════════════════════════════

\begin{document}

% ── Headline block ───────────────────────────────────────────────────────────
\begin{center}
  {\LARGE\bfseries\color{SectionBlue}
    vox Gives AI Coding Agents a Voice and Audio Layer ---\\
    Developers Work Alongside Their Agent by Ear, Not Just Eyes} \\[0.8em]
  {\large\color{AccentGray}
    vox is the voice and audio layer for AI coding agents: voice
    personalities, mood-aware expressive delivery, background music, and
    spoken summaries through five TTS providers --- including zero-config
    offline fallbacks. Because listening and watching the screen draw on
    largely separate cognitive channels, developers track their agent's
    progress by ear while their eyes stay on other work.}
\end{center}

\vspace{0.5em}

% ══════════════════════════════════════════════════════════════════════════════
% PRESS RELEASE
% ══════════════════════════════════════════════════════════════════════════════

\section*{Press Release}

SAN FRANCISCO --- September 2026 --- Today, punt-labs announced the general
availability of vox, a free, open-source voice and audio layer for AI
coding agents. It gives an agent a voice --- distinct voices from a roster,
mood-aware expressive delivery, background music, and quips --- so working
alongside Claude Code feels less like watching a console scroll and more like
collaborating with something that has a voice. Because listening and watching
the screen draw on largely separate cognitive resources
\cite{baddeleyhitch1974workingmemory,baddeley2000episodicbuffer}, a developer
can track the agent's progress by ear without keeping their eyes on a wall of
console text \cite{springer2016audiovisual,mccrickardchewar2003attuning}.
Audio is delivered through premium
text-to-speech from ElevenLabs, OpenAI, or AWS Polly, with zero-config offline
fallbacks via macOS \texttt{say} and Linux \texttt{espeak-ng}, a one-command
install, and graceful absence when not installed. For developers with
multi-machine setups, vox supports remote audio: run Claude Code on a
headless server and hear it on the local workstation
(see \faqref{faq:competitors}).

\prsection{Problem}

Working with an AI coding agent is, today, a visual-only experience. The
agent's reasoning, its progress, and its output all arrive as a wall of console
text the developer has to keep watching to stay aware of --- so staying engaged
means staying glued to the terminal, and looking away means losing the thread.
So developers multitask. They start a refactor, switch to a document, answer a
Slack thread. The agent keeps working --- for seconds, sometimes for thirty
minutes or more on autonomous tasks \cite{anthropic2026longrunnagents}. When it finishes, or when it needs
permission to proceed, the developer does not know until they check the
terminal. On complex tasks, the agent stops to ask for clarification more than
twice as often as the developer interrupts it
\cite{anthropic2026agentautonomy}. While auto-approve rates are growing ---
from 20\% for new users to over 40\% for experienced users --- the majority
of sessions still involve at least one approval gate, and each unnoticed
prompt is dead time: the agent waits, the developer works on something else,
and neither makes progress on the original task.

The cost is cumulative. Reactive interruptions cost over 23 minutes of
recovery time on average \cite{mark2008interrupted}. Self-initiated checks
--- glancing at the terminal to see if the agent is done --- are lower-cost
than reactive interruptions, but each check still breaks the developer's
current focus. Across a day of agentic work, these small breaks add up:
51\% of professional developers now use AI tools daily
\cite{stackoverflow2025survey}, and a growing fraction use autonomous agents
that run for extended periods
(see \faqref{faq:tam}). And the trend is accelerating: IDC projects 80\% of
enterprise apps will embed AI copilots by end of 2026
\cite{couriernotificationfatigue2026}, so the volume of agent activity
competing for a developer's attention will only grow.

\prsection{Solution}

vox adds an audio layer to Claude Code through slash commands. Together
they do two things: they give the session a voice and a mood, and they let the
developer stay aware of the agent by ear. Notifications are the first of these
uses, not the only one.

\texttt{/vox y} enables chime notifications for two events: task completion
and permission prompts. When the agent finishes a multi-file refactor, the
developer hears an audio tone. When the agent needs approval to run a
destructive command, the developer hears that too --- no more silent permission
dialogs burning minutes of idle time. \texttt{/vox c} adds continuous mode
with spoken summaries: milestone updates during long-running tasks, so the
developer hears progress without checking in. \texttt{/unmute} enables voice
mode with a specific voice from the roster.

\texttt{/mute} switches from voice to chime --- the same events fire, but
the notification is an audio tone instead of words. Chimes are mood-aware:
when a vibe is active, tones pitch-shift to match the session mood (brighter
when things are going well, darker when they're not). Eight distinct signals
across three mood variants keep chime-only mode as expressive as voice. For
developers who want awareness without narration.

\texttt{/vibe} sets the session mood, and the voice delivers to match it ---
brighter and more energetic during a green test run, weary after a string of
failures, with expressive cues (\texttt{[excited]}, \texttt{[weary]},
\texttt{[sighs]}) woven into premium-provider speech. The point is not
character role-play but a session that sounds like it is going the way it is
going --- a collaborator with a voice, not a status bar.

\texttt{/recap} is on-demand: after a wall of terminal output, the developer
types \texttt{/recap} and hears a 30-second spoken summary of what just
happened. No scrolling, no parsing diffs. The agent extracts the key points
and speaks them while the developer scans the code
(see \faqref{faq:getting-started}).

\texttt{/music on} starts vibe-driven background music that loops during coding
sessions. The music adapts when the session mood changes: upbeat during green
test runs, subdued during debugging. Music plays at reduced volume under speech
and chimes. vox is partly entertainment by design here: the agent plays
DJ, curating vibe-matched background music pools from a music panel with a
DJ-booth personality.

The plugin auto-detects the best available TTS provider --- ElevenLabs for
natural voice quality, OpenAI for low latency, AWS Polly for cost efficiency,
macOS \texttt{say} and Linux \texttt{espeak-ng} for zero-config offline use
--- and requires no voice selection or audio device configuration. For
developers who SSH into remote machines, setting \texttt{VOXD\_HOST},
\texttt{VOXD\_PORT}, and \texttt{VOXD\_TOKEN} routes audio to a local
\texttt{voxd} daemon so
notifications play on the workstation, not the server
(see \faqref{faq:competitors}).

\prsection{Customer Quote}

\begin{customerquote}
\textit{``I run Claude Code on a second monitor for refactoring work. Before
vox, I'd context-switch back every few minutes to check if it was done or
stuck on a permission prompt. The worst was realizing it had been waiting for
approval for twenty minutes while I was in a design review. Now I hear `task
complete' or `needs your approval' and I know exactly when to switch back.
Last Thursday I went through a full hour-long code review without once checking
the terminal --- and I didn't miss a single prompt. And it's more than the
alerts: when it's grinding through a migration I can hear it working, the mood
shifts when the tests go green, and the session stopped being a silent wall of
text I had to babysit.''}

\raggedleft --- \textbf{Marcus Chen}, Senior Software Engineer at a Series B
startup
\end{customerquote}

\prsection{Getting Started}

One command installs everything:

{\small\texttt{curl -fsSL https://raw.githubusercontent.com/punt-labs/vox/main/install.sh | sh}}

The script checks Python 3.13+, installs \texttt{uv} if missing, installs
the CLI, registers the MCP server with Claude Code, and runs
\texttt{vox doctor} to verify that at least one TTS provider is available.
Restart Claude Code and type \texttt{/vox y}. From that moment, every
task completion and permission prompt produces a chime notification.
Type \texttt{/vox c} for continuous spoken summaries.

There is no voice selection wizard and no audio device configuration. The
\texttt{voxd} daemon handles synthesis, caching, and playback as a system
service. If \texttt{vox doctor} passes, everything works. If vox is not
installed, Claude Code behaves exactly as before --- no error, no
degradation.

\prsection{Spokesperson Quote}

\begin{spokespersonquote}
``We built vox because we kept walking away from the terminal and losing
time. The agent finishes, or worse, it sits waiting for permission, and nobody
knows. Voice solves this in a way that visual notifications cannot ---
it reaches you when your eyes are on another screen. But notifications were
just the start. Listening and watching the screen pull on different parts of
your attention, so a spoken update lands while your eyes stay on the code --- and a
voice with a mood makes a long session feel less like babysitting a log and
more like working next to someone. We chose premium TTS
providers because the OS built-in voice is unpleasant enough that people turn
it off, which defeats the purpose. If the notification sounds good, people
leave it on. That is the entire design philosophy: voice that supplements text,
never replaces it, and sounds good enough to keep.''

\raggedleft --- \textbf{JM Freeman}, Creator of vox
\end{spokespersonquote}

\prsection{Call to Action}

vox is available now at
\url{https://github.com/punt-labs/vox}. It is free, open source (MIT license),
and always will be. Install with one command:

{\small\texttt{curl -fsSL https://raw.githubusercontent.com/punt-labs/vox/main/install.sh | sh}}

Restart Claude Code and type \texttt{/vox y}.

\newpage

% ══════════════════════════════════════════════════════════════════════════════
% FREQUENTLY ASKED QUESTIONS
% ══════════════════════════════════════════════════════════════════════════════

\section*{Frequently Asked Questions}

% ── External FAQs (Customer-facing) ──────────────────────────────────────────

\subsection*{External FAQs}

\begin{faqpair}{What is vox and who is it for?}\label{faq:what}
  vox is the voice and audio layer for AI coding agents. It gives an
  agent a voice --- distinct voices from a roster, mood-aware expressive
  delivery, background music, and quips --- so working with Claude Code feels
  less like watching a console and more like collaborating with something that
  has a voice. Because listening and watching the screen draw on largely
  separate cognitive resources
  \cite{baddeleyhitch1974workingmemory,baddeley2000episodicbuffer}, it also
  lets a developer track the agent's progress by ear without watching the
  terminal \cite{springer2016audiovisual,mccrickardchewar2003attuning}.
  It is for
  developers who use Claude Code for autonomous tasks --- refactoring,
  debugging, code generation --- and want the session to be both more engaging
  and easier to stay aware of while they work on something else. This includes
  developers who run Claude Code on remote servers via SSH: remote audio
  routes the sound to the local workstation where the developer can hear it.
\end{faqpair}

\begin{faqpair}{How is this different from agent-notify or AgentVibes?}\label{faq:competitors}
  Three open-source tools have added voice to coding agents. \textbf{AgentVibes}
  \cite{agentvibes2025} is the strongest --- an actively maintained, feature-rich
  Claude Code companion. \textbf{agent-notify} \cite{agentnotify2025} is minimal:
  OS \texttt{say} plus push alerts. A third, claude-code-voice-handler, used
  OpenAI TTS but appears to have been taken down. The real comparison is with
  AgentVibes, so here is an honest head-to-head.

  \begin{center}
  \small
  \renewcommand{\arraystretch}{1.3}
  \begin{tabularx}{\linewidth}{@{} l X X @{}}
  \toprule
  \textbf{\color{SectionBlue}Dimension} & \textbf{\color{SectionBlue}AgentVibes} & \textbf{\color{SectionBlue}vox} \\
  \midrule
  Providers &
    Six, including ElevenLabs (added v5.11.0, 2026-06-27) and local neural
    Piper and Kokoro \cite{agentvibes2026releasenotes,agentvibes2026providers} &
    ElevenLabs, OpenAI, Polly, plus offline \texttt{say} / \texttt{espeak-ng} \\
  Expressive delivery &
    19 static, command-invoked personalities (pirate, zen, \dots); no automatic
    mood \cite{agentvibes2026personalities} &
    \texttt{/vibe} auto-mode tracks the session and shifts delivery
    ([excited], [weary], [sighs]) \\
  Music &
    Bundled MP3 catalog with per-track volume &
    Generated 12-track vibe-matched pools via the ElevenLabs Music API \\
  Surfaces &
    CLI (\texttt{npx agentvibes}) + MCP server, both wrapping bash hook scripts
    \cite{agentvibes2026technicaldeepdive} &
    MCP server + importable Python library (\texttt{VoxClient}) + CLI, one daemon \\
  Reusable engine &
    ``Primarily a Claude Code plugin system,'' not a reusable library
    \cite{agentvibes2026technicaldeepdive} &
    langlearn-tts embeds \texttt{VoxClient} as its speech backend \\
  Stack &
    Roughly 65\% JavaScript; \texttt{blessed} multi-tab TUI
    \cite{agentvibes2026github} &
    Python daemon plus thin stdlib clients \\
  Age &
    Repo created 2025-06-10 \cite{agentvibes2026github} &
    First commit 2026-02-25 (\textasciitilde{}8.5 months later) \\
  \bottomrule
  \end{tabularx}
  \renewcommand{\arraystretch}{1.0}
  \end{center}

  \textbf{Where vox is genuinely different.} Not voice quality --- AgentVibes
  now ships ElevenLabs too, so both reach the same premium tier. vox's edge is
  structural. It exposes an \emph{importable engine}: a downstream application
  can embed \texttt{VoxClient}/\texttt{VoxClientSync} in-process, which
  langlearn-tts already does. AgentVibes, by its own technical documentation, is
  ``primarily a Claude Code plugin system rather than a universally reusable
  library'' \cite{agentvibes2026technicaldeepdive} --- its MCP server wraps bash
  hook scripts rather than exposing an embeddable SDK. AgentVibes does have a
  working CLI and MCP server, so this is a library-reusability edge, not a claim
  that it lacks interfaces. Beyond that: vox's mood is dynamic and automatic
  where AgentVibes' personalities are static and manually selected; vox's music
  is generated and vibe-matched where AgentVibes plays a fixed catalog; and vox
  is a focused Python engine rather than a JavaScript application with a full
  multi-tab TUI.

  \textbf{Where AgentVibes leads, and what vox can learn.} AgentVibes is the
  older, more established project and bundles local neural voice models ---
  Piper (900+ voices) and Kokoro (non-English, CPU-only)
  \cite{agentvibes2026providers} --- giving it a far broader \emph{offline}
  catalog than vox, whose offline fallback is only macOS \texttt{say} and Linux
  \texttt{espeak-ng}. That is a real gap: a bundled local neural model is the one
  concrete idea worth taking from AgentVibes if vox wants a stronger offline
  story.
\end{faqpair}

\begin{faqpair}{How do I get started?}\label{faq:getting-started}
  One command:

  {\small\texttt{curl -fsSL https://raw.githubusercontent.com/punt-labs/vox/main/install.sh | sh}}

  The script installs the CLI, registers the MCP server with Claude Code, and
  runs \texttt{vox doctor} to verify providers. Restart Claude Code. Type
  \texttt{/vox y}. Notifications are now active. Type \texttt{/recap} after
  any long output for a spoken summary.

  For manual installation:
  {\small\texttt{uv tool install punt-vox \&\& vox install \&\& vox doctor}}.
\end{faqpair}

\begin{faqpair}{How much does it cost?}\label{faq:cost}
  vox itself is free (MIT license). TTS provider API calls have costs:

  \begin{itemize}[leftmargin=1.5em, nosep]
    \item ElevenLabs: from \$5/month for 30{,}000 characters
          \cite{texttolab2025pricing}
    \item OpenAI TTS: \$15 per million characters (tts-1)
          \cite{texttolab2025pricing}
    \item AWS Polly: \$4.80 per million characters (Standard)
          \cite{texttolab2025pricing}
  \end{itemize}

  A typical notification is 50--200 characters. At OpenAI's rate, 1{,}000
  notifications cost approximately \$0.003. Even heavy daily use (50
  notifications/day) runs under \$0.10/month. The dominant cost is the API
  key setup, not the usage.
\end{faqpair}

\begin{faqpair}{What happens to my data?}\label{faq:data}
  vox sends the text of each notification to your chosen TTS provider for
  synthesis. The provider receives the text, returns audio, and vox plays
  it locally. No data is stored by vox beyond a local audio cache in
  \texttt{\texttildelow/.punt-labs/vox/cache/} and saved audio in
  \texttt{\texttildelow/Music/vox/} (configurable via
  \texttt{VOX\_OUTPUT\_DIR}). vox has no telemetry, no analytics, and no
  cloud service of its own.

  The privacy implication: your TTS provider sees the text of your
  notifications. For ElevenLabs and OpenAI, review their data retention
  policies. For AWS Polly, standard AWS data handling applies. If this is
  unacceptable, vox is not the right tool --- the existing OS TTS
  alternatives process text locally with no external API calls.
\end{faqpair}

\begin{faqpair}{How can I use vox --- MCP server, Python library, or CLI?}
  The Claude Code plugin (\texttt{/vox}, \texttt{/unmute}, \texttt{/mute},
  \texttt{/recap}, \texttt{/vibe}, \texttt{/music}) is one way to use vox, not
  the only way. vox ships three first-class surfaces, and a reader should come
  away seeing that it is usable well outside a Claude Code plugin:

  \textbf{MCP server.} \texttt{vox-server} exposes vox to any
  MCP-compatible agent under the key \texttt{mic}. The agent calls a tool ---
  for example \texttt{unmute} or \texttt{speak} --- to make itself heard. This
  is the same surface the Claude Code plugin drives under the hood.

  \textbf{Python library.} \texttt{VoxClient} and \texttt{VoxClientSync}
  (from \texttt{vox.client}) give any Python program synthesis, recording,
  voice listing, health checks, and music-program control:

  {\small\begin{verbatim}
async with VoxClient() as vox:
    await vox.synthesize("Build finished.")
\end{verbatim}}

  langlearn-tts already uses this surface as its speech backend.

  \textbf{CLI.} The \texttt{vox} command drives everything from a shell ---
  \texttt{say}, \texttt{record}, \texttt{voice}, \texttt{voices},
  \texttt{vibe}, \texttt{notify}, \texttt{speak}, \texttt{music},
  \texttt{daemon}, \texttt{status}, and \texttt{doctor}:

  {\small\begin{verbatim}
vox say "Build finished."
\end{verbatim}}

  All three surfaces talk to the same \texttt{voxd} daemon, so voice, caching,
  and playback behave identically no matter which you use. Integrating with
  Cursor, GitHub Copilot, or another agent means wiring one of these surfaces
  into that tool's hook or notification system.
\end{faqpair}

\begin{faqpair}{Can I use my own voice or choose a specific voice?}
  vox auto-selects a default voice per provider (Matilda for ElevenLabs,
  Nova for OpenAI, Joanna for Polly). Use \texttt{/unmute @matilda} to set a
  session voice, or \texttt{/unmute @} to browse the roster. The CLI accepts
  \texttt{-{}-voice} flags for overriding the default. The design philosophy
  is zero configuration --- defaults work out of the box, but voice selection
  is available for those who want it.
\end{faqpair}

% ── Internal FAQs (Business-facing) ──────────────────────────────────────────

\newpage
\subsection*{Internal FAQs}

\subsubsection*{Value \& Market}

\begin{faqpair}{What is the total addressable market?}\label{faq:tam}
  Bottoms-up estimate from three data points:

  \textbf{AI coding agent users:} 31\% of professional developers use AI
  agents \cite{stackoverflow2025survey}. GitHub Copilot alone has 20 million
  cumulative users \cite{techcrunch2025copilot20m}. Cursor has over 1 million
  users \cite{opsera2025cursor}. Claude Code is used by 41\% of AI agent users
  \cite{stackoverflow2025survey}.

  \textbf{Autonomous task runners:} Not all agent users run autonomous tasks.
  Anthropic data shows the median Claude Code turn is approximately 45 seconds,
  but the 99.9th percentile grew from under 25 minutes to over 45 minutes
  between October 2025 and January 2026 \cite{anthropic2026agentautonomy}. The
  growth is driven by user behavior --- granting more autonomy --- not just
  model capability. Background tasks are explicitly recommended for work
  exceeding 30 seconds \cite{venturebeat2025claudetasks}.

  \textbf{Target segment:} Developers who run autonomous agent tasks and
  multitask during execution. Applying the math: if roughly 28 million
  professional developers exist globally (assumption; order-of-magnitude,
  uncited), 31\% use agents (~8.7M), 41\%
  of agent users use Claude Code (~3.6M), and 10--20\% run autonomous tasks
  long enough to warrant notification, the target segment is 360{,}000 to
  720{,}000 developers. This is the upper bound; actual reachable users via
  PyPI are a fraction of this.

  Because vox is free, TAM is measured in adoption, not revenue. Success
  means becoming the default voice layer that Claude Code users install
  alongside the agent.
\end{faqpair}

\begin{faqpair}{What evidence do we have that developers want this?}\label{faq:evidence}
  Three categories of evidence, all indirect. No primary customer data exists
  yet (hypothesis stage).

  \textbf{Behavioral:} Two open-source tools already exist to solve this
  problem. agent-notify \cite{agentnotify2025} adds sound and voice
  notifications to Claude Code, Cursor, and Codex. AgentVibes
  \cite{agentvibes2025} adds voice, background music, and personality styles
  to Claude Code. The existence of these tools --- built by independent
  developers, not companies --- is evidence that the pain point is real enough
  for people to build solutions.

  \textbf{Structural:} Anthropic's own data confirms the structural
  conditions. Agents ask for clarification more than twice as often as users
  interrupt on complex tasks \cite{anthropic2026agentautonomy}. The hooks
  system exposes \texttt{Notification} events with
  \texttt{permission\_prompt} and \texttt{idle\_prompt} subtypes
  \cite{claudecodehooks2026} --- Anthropic built the infrastructure for
  exactly this kind of notification.

  \textbf{Adjacent:} The channel-separation premise under vox is
  well grounded. In Baddeley's working-memory model the phonological loop
  (verbal/acoustic) and the visuospatial sketchpad (visual/spatial) are
  separable, capacity-limited stores
  \cite{baddeleyhitch1974workingmemory,baddeley2000episodicbuffer}, so
  attending to audio and attending to on-screen text draw on largely separate
  resources rather than competing for one. HCI notification research shows
  audio notifications maintain workspace awareness without demanding visual
  attention \cite{springer2016audiovisual}, and peripheral/audio channels
  lower the interruption cost of staying aware relative to a visual channel
  that requires focus \cite{mccrickardchewar2003attuning}. This supports two
  narrow claims --- that attending to audio and attending to on-screen text
  draw on largely separate channels, and that audio lets a developer track
  progress without visual attention --- and nothing stronger: we do not claim a spoken summary carries
  less total cognitive load than reading the underlying output, which the
  literature does not establish.

  \textbf{What is missing:} No user interviews. No usage data. No measurement
  of how long permission prompts sit unnoticed. This is the most important
  research gap
  (see \faqref{faq:validation-plan}).
\end{faqpair}

\begin{faqpair}{Who are the competitors and why will we win?}\label{faq:competitive}
  \textbf{agent-notify} \cite{agentnotify2025}: Open-source. Sound alerts,
  macOS \texttt{say}, ntfy push notifications. Supports Claude Code, Cursor,
  Codex. Strengths: no API key required, cross-agent support. Weakness: OS
  TTS only, no packaging or test suite, single developer.

  \textbf{AgentVibes} \cite{agentvibes2025}: Open-source, actively maintained,
  and the older project (repo created 2025-06-10). Six TTS providers, including
  ElevenLabs since v5.11.0 (2026-06-27) plus local neural Piper (900+ voices)
  and Kokoro \cite{agentvibes2026releasenotes,agentvibes2026providers}. Bundled
  music catalog, 19 static personality styles, verbosity levels. Strengths: most
  feature-complete, broad offline voice catalog, creative UX. Weakness: no
  importable/reusable TTS library (the MCP server wraps bash hook scripts)
  \cite{agentvibes2026technicaldeepdive}, static personalities rather than
  dynamic mood, no PyPI packaging or test suite.

  \textbf{OS notification center:} macOS/Linux system notifications. Always
  available, no setup. Weakness: visual-only, competes with every other app's
  notifications, easy to miss during focus work.

  \textbf{Why vox wins:} Voice quality anchors a broader bet. The macOS
  \texttt{say} voice is recognizable and unpleasant enough that users turn it
  off. ElevenLabs consistently ranks among the top TTS providers in blind
  listening tests and has \$330M+ ARR \cite{cnbc2026elevenlabs}. If the audio
  sounds good, people leave it on --- and what they leave on is more than
  notifications: voice personalities, mood-aware expressive delivery, and
  music make the session engaging, while eyes-free progress tracking (audio
  and on-screen visual work use largely separate channels) keeps it useful
  during focus work. This is the core bet: an engaging, genuinely usable audio layer,
  anchored by premium voice quality, drives sustained adoption where OS TTS
  does not.

  \textbf{Honest assessment of risk:} AgentVibes added ElevenLabs in v5.11.0
  (2026-06-27) \cite{agentvibes2026releasenotes}, so the two tools now reach
  provider parity --- premium voice quality is no longer a vox differentiator.
  vox's remaining edge is structural, not tonal: a reusable, importable engine
  (\texttt{VoxClient}, already embedded by langlearn-tts) where AgentVibes is
  ``primarily a Claude Code plugin system''
  \cite{agentvibes2026technicaldeepdive}; dynamic, automatic mood versus 19
  static personalities; generated, vibe-matched music versus a bundled catalog;
  and a focused Python engine versus a roughly 65\% JavaScript application with
  a multi-tab TUI \cite{agentvibes2026github}. AgentVibes leads on offline voice
  breadth (bundled Piper and Kokoro models) and is the older, more established
  project. The competitive moat is narrow and rests on reusability and focus,
  not voice.
\end{faqpair}

\begin{faqpair}{What is the user acquisition strategy?}
  Three channels:

  \textbf{PyPI discoverability:} Users searching for ``claude code voice'' or
  ``tts mcp'' find vox. The install script auto-configures MCP, reducing
  setup to one command.

  \textbf{Claude Code plugin marketplace:} When the marketplace launches,
  vox will be listed as a plugin. One-click install with
  \texttt{/install punt-vox}.

  \textbf{Word of mouth:} Voice notifications are inherently demonstrable.
  When a developer's agent speaks ``task complete'' during a pairing session or
  meeting, the person next to them asks what that was. This is the strongest
  acquisition vector for an audio product.

  No paid acquisition planned. The product is free; the audience is technical;
  organic distribution is the appropriate channel.
\end{faqpair}

\begin{faqpair}{What is your next step to validate your vision?}\label{faq:validation-plan}
  All core features are shipped (v4.12.5). The product is approaching
  public recommendation. Validation now focuses on adoption and retention:

  \begin{enumerate}[leftmargin=1.5em, nosep]
    \item Do users who enable \texttt{/vox y} leave it on after one week?
    \item Which event fires more often: task completion or permission prompt?
    \item Do users stay on chimes (\texttt{/vox y}) or upgrade to voice
          (\texttt{/unmute})?
    \item Does remote \texttt{voxd} expand the addressable audience to
          SSH-based workflows?
  \end{enumerate}

  If retention after one week is below 30\%, the value proposition is
  wrong and we should stop. If permission-prompt notifications dominate,
  that reframes the product from ``voice for your agent'' to ``never miss
  a permission prompt again.'' Both outcomes inform whether to continue.

  Target: 5 developers using vox daily for 2 weeks. Collect
  feedback via GitHub Discussions.
\end{faqpair}

\subsubsection*{Technical}

\begin{faqpair}{What are the major technical risks?}\label{faq:tech-risks}
  \textbf{Claude Code hooks stability (medium risk):} vox depends on
  Claude Code's hooks system --- specifically the \texttt{Stop} event and
  \texttt{Notification} event with \texttt{permission\_prompt} subtype
  \cite{claudecodehooks2026}. Hooks are documented first-party API, but
  breaking changes could disable notifications. Mitigation: hooks are a
  foundational feature of Claude Code's extensibility model, used by many
  plugins. Breaking them would affect the entire plugin ecosystem.

  \textbf{TTS provider latency (low risk):} Voice notifications must arrive
  within 1--2 seconds of the triggering event to feel real-time. ElevenLabs
  streaming API and OpenAI TTS both return first audio within 200--500ms.
  AWS Polly is faster still. Mitigation: auto-detect the lowest-latency
  available provider for notifications; reserve higher-quality providers for
  \texttt{/recap} summaries where latency is less critical.

  \textbf{Audio playback on headless/remote systems (resolved, shipped
  v4.7.7):} Developers using Claude Code via SSH or in cloud VMs previously
  had no audio path. As of v4.7.7, \texttt{VOXD\_HOST}/\allowbreak\texttt{VOXD\_PORT}/\allowbreak\texttt{VOXD\_TOKEN}
  env vars route synthesis requests to a remote \texttt{voxd} instance
  running on the developer's local workstation. Token-based authentication
  secures the connection. \texttt{vox doctor} reports active env var
  overrides. For local-only setups, \texttt{vox doctor} still checks audio
  playback capability at install time and notifications fail silently when
  audio is unavailable.

  All core technical risks have been resolved. The TTS engine (synthesis,
  batching, merging) is built and tested. The Stop hook uses a decision-block
  pattern where the model generates a 1-2 sentence summary and calls the TTS
  tool. Vibe tags are resolved deterministically from accumulated session
  signals, written to a gitignored ephemeral config file
  (\texttt{vox.local.md}), and picked up automatically by the TTS tool ---
  no data passes through user-visible output.
\end{faqpair}

\begin{faqpair}{What dependencies exist on external systems?}
  \textbf{Claude Code hooks:} Required for \texttt{/vox} notifications to fire
  automatically. Without hooks, vox still works as a manual TTS tool
  (\texttt{/recap}, CLI synthesis) but loses its primary value proposition.

  \textbf{TTS provider APIs:} At least one of ElevenLabs, OpenAI, or AWS
  Polly must be reachable. Auto-detection falls back through the chain.
  If all providers are unavailable, \texttt{vox doctor} reports the failure
  and no audio plays.

  \textbf{pydub + ffmpeg:} Used for audio merging and format conversion.
  ffmpeg is required at runtime. \texttt{vox doctor} checks for its presence.

  \textbf{MCP protocol:} The plugin registers as an MCP server. MCP is
  maintained by Anthropic and adopted by multiple clients. Low risk of
  abandonment.

  No runtime cloud dependencies exist beyond the TTS API calls themselves.
  vox has no backend service, no database, and no telemetry. The optional
  remote \texttt{voxd} mode connects to a user-controlled daemon on another
  machine --- not to any punt-labs service.
\end{faqpair}

\begin{faqpair}{What is the estimated development timeline?}\label{faq:timeline}
  punt-vox is at v4.12.5 on PyPI (July 2026). The product is nearly
  feature-complete.

  \textbf{Shipped:} Stop hook notification with spoken summaries.
  \texttt{/vox y|n|c} toggle. \texttt{/unmute}/\texttt{/mute} for voice
  vs.\ chime. Permission-blocked notification via \texttt{Notification}
  hook. \texttt{/recap} command (agent-executed spoken summary).
  \texttt{/vibe} with auto-detection and manual modes --- auto-vibe uses
  deterministic signal-to-tag mapping (no LLM needed) to resolve
  ElevenLabs expressive tags from accumulated session signals.
  Five providers (ElevenLabs, OpenAI, Polly, macOS \texttt{say}, Linux
  \texttt{espeak-ng}). Mood-aware chimes with pitch shifting
  (8 signals $\times$ 3 moods). Non-blocking MCP tools.
  \texttt{voxd} audio daemon with WebSocket protocol (v3.0.0) ---
  separates synthesis/playback from MCP server for fast hook dispatch
  (observed on the order of a few milliseconds; not formally benchmarked).
  Background music generation via ElevenLabs Music API (v4.4.0).
  Per-call API key passthrough for billing isolation (v4.3.0).
  Config split into tracked \texttt{vox.md} and ephemeral
  \texttt{vox.local.md} (v4.7.5).
  Claude Code plugin on marketplace. CLI global flags
  (\texttt{-{}-json}, \texttt{-{}-verbose}, \texttt{-{}-quiet}).
  Remote \texttt{voxd} connectivity ---
  \texttt{VOXD\_HOST}/\texttt{VOXD\_PORT}/\texttt{VOXD\_TOKEN} env vars
  for cross-host audio playback without SSH tunnels, plus server-side
  \texttt{VOXD\_BIND} to control the bind address (v4.7.7).

  \textbf{Next:} Downstream library adoption beyond langlearn-tts. Public
  recommendation and outreach.
\end{faqpair}

\begin{faqpair}{What is the scaling story?}
  vox runs on the developer's machine. There is no server to scale.
  Each notification is an independent TTS API call --- no state accumulates,
  no queues build up.

  At the TTS provider level: ElevenLabs handles concurrent requests across
  millions of users (\$330M+ ARR, 41\% Fortune 500 adoption
  \cite{cnbc2026elevenlabs}). OpenAI and AWS Polly are similarly robust.
  vox will never be a meaningful fraction of any provider's load.

  The only local scaling concern is audio file accumulation. Saved audio
  lands in \texttt{\texttildelow/Music/vox/} (configurable); the synthesis
  cache in \texttt{\texttildelow/.punt-labs/vox/cache/} is content-addressed
  and can be cleared via \texttt{vox cache clear}.
\end{faqpair}

\subsubsection*{Business}

\begin{faqpair}{What is the revenue model?}
  There is none. vox is free, open-source software under the MIT license.
  The project exists to solve a real problem in the Claude Code ecosystem and
  to serve as the TTS engine for other punt-labs tools (langlearn-tts uses
  punt-vox as its speech backend). Value is measured in adoption and ecosystem
  contribution, not revenue.
\end{faqpair}

\begin{faqpair}{What does maintaining vox cost?}
  \textbf{Infrastructure:} Zero. vox runs on the user's machine. GitHub
  Actions CI is free for public repos.

  \textbf{TTS API costs during development:} The test suite mocks all
  providers, so CI incurs zero TTS costs. Manual integration testing against
  live providers costs under \$0.10 per session. Negligible.

  \textbf{Maintainer time:} The primary cost. Estimated at 2--5 hours per
  week during active development (Phases 1--3), dropping to 1--2 hours per
  week for maintenance (bug fixes, dependency updates, community triage).

  \textbf{Opportunity cost:} Maintainer time spent on vox is time not
  spent on other punt-labs projects. If adoption metrics do not meet
  thresholds after 3 months (see \faqref{faq:metrics}), maintenance mode is
  the appropriate response --- fix bugs, accept PRs, but no new features.
\end{faqpair}

\begin{faqpair}{What are the key metrics for success?}\label{faq:metrics}
  \begin{itemize}[leftmargin=1.5em, nosep]
    \item \textbf{PyPI installs:} Monthly download count as adoption proxy.
          Target: 500 monthly installs within 6 months of the notification
          feature shipping.
    \item \textbf{Notification retention:} Percentage of users who enable
          \texttt{/vox y} and keep it enabled after one week. Target:
          30\%+. Below this, the product is not sticky.
    \item \textbf{Event distribution:} Ratio of task-completion to
          permission-prompt notifications. Informs whether the product's
          primary value is ``hear when done'' or ``never miss a prompt.''
    \item \textbf{GitHub stars:} Community interest signal. Target: 200
          stars within 6 months.
    \item \textbf{Downstream dependents:} Other tools using punt-vox as a
          library. Current: 1 (langlearn-tts). Target: 3+ within a year.
  \end{itemize}

  Decision threshold: if notification retention is below 30\% after the
  initial user cohort, reconsider the product direction before investing in
  Phase 3 features.
\end{faqpair}

\begin{faqpair}{Why now? What has changed?}
  Three converging shifts:

  \textbf{Autonomous agents are real.} Claude Code's 99.9th percentile task
  duration grew from under 25 minutes to over 45 minutes in three months
  \cite{anthropic2026agentautonomy}. Background tasks are a documented
  workflow \cite{venturebeat2025claudetasks}. Agents now run long enough that
  developers routinely leave them unattended.

  \textbf{Hooks enable it.} Claude Code's hooks system exposes lifecycle
  events --- \texttt{Stop}, \texttt{Notification}, \texttt{TaskCompleted}
  --- as structured, documented API \cite{claudecodehooks2026}. Before hooks,
  notification plugins required polling or terminal scraping. Hooks make
  event-driven notifications clean and reliable.

  \textbf{Premium TTS is accessible.} ElevenLabs, OpenAI, and Polly all offer
  API-based synthesis at costs below \$0.01 per notification. Two years ago,
  natural-sounding TTS required expensive enterprise contracts. Today it costs
  less than a fraction of a cent per utterance
  \cite{texttolab2025pricing}.
\end{faqpair}

\begin{faqpair}{What are we not building?}
  Covered in the Feature Appendix (Won't Do section). The boundaries that
  matter most:

  \begin{itemize}[leftmargin=1.5em, nosep]
    \item \textbf{Not voice input.} vox speaks to the developer. It does
          not listen. Voice commands for agents are a different product.
    \item \textbf{Not a screen reader.} Accessibility narration of terminal
          output requires fundamentally different architecture (continuous
          reading, cursor tracking, ARIA-like semantics). vox is
          event-driven notifications, not continuous narration.
  \end{itemize}
\end{faqpair}

\begin{faqpair}{What does failure look like?}\label{faq:failure}
  Two plausible failure scenarios:

  \textbf{People turn it off.} Voice notifications are intrusive. If the
  default voice is too loud, too frequent, or too slow, users disable
  \texttt{/vox n} within hours and never re-enable. The notification
  retention metric (see \faqref{faq:metrics}) is the canary: below 30\%
  after one week means the UX is wrong.
  \textit{Response:} \texttt{/mute} (chime-only) is the escape valve.
  If retention is low on voice but acceptable on chime, the product is
  an audio notification tool, not a voice tool. Adjust positioning
  accordingly.

  \textbf{The pain point shrinks.} Anthropic's data shows auto-approve rates
  growing from 20\% to 40\%+ as users gain experience
  \cite{anthropic2026agentautonomy}. If the trend continues, permission
  prompts become rare and the strongest use case
  (``never miss a prompt'') weakens. At the same time, longer autonomous
  runs make the task-completion notification more valuable.
  \textit{Response:} Monitor the event distribution metric. If permission
  prompts decline to less than 10\% of notifications, pivot messaging to
  ``hear when your agent finishes'' and invest in \texttt{/recap} quality.
\end{faqpair}

% ══════════════════════════════════════════════════════════════════════════════
% FOUR RISKS ASSESSMENT
% ══════════════════════════════════════════════════════════════════════════════

\newpage
\section*{Risk Assessment}

\renewcommand{\arraystretch}{1.6}
\begin{tabularx}{\textwidth}{@{} l l X @{}}
\toprule
\textbf{\color{SectionBlue}Risk} & \textbf{\color{SectionBlue}Rating} & \textbf{\color{SectionBlue}Assessment} \\
\midrule
Value      & \textcolor{RiskAmber}{\textbf{Medium}} &
  Demand is proven by the existence of two competing tools. The bet is broader
  than notifications: that an engaging audio layer --- voice personalities,
  mood-aware delivery, music --- plus eyes-free progress tracking (audio and
  on-screen visual work draw on largely separate channels) drives sustained
  adoption, anchored by premium voice quality that OS TTS cannot match. No
  primary user data yet. This broadened bet --- engagement, mood, music ---
  carries more unvalidated surface than the notification claim it replaced, so
  Medium is a floor, not a settled read, until the validation plan produces
  usage data. \\
Usability  & \textcolor{RiskGreen}{\textbf{Low}} &
  One-command install and auto-detection of providers. Since v0.7.0, offline
  fallbacks (macOS \texttt{say}, Linux \texttt{espeak-ng}) work with zero
  API keys --- install and go. Premium providers still require account
  creation and environment variable configuration, but the onboarding path
  no longer requires it. \\
Feasibility & \textcolor{RiskGreen}{\textbf{Low}} &
  The TTS engine, hooks integration, \texttt{voxd} daemon, background music,
  config split, remote audio, and all commands are built and shipped.
  No open design questions remain. \\
Viability  & \textcolor{RiskGreen}{\textbf{Low}} &
  No revenue model to validate. Zero infrastructure cost. MIT license.
  The only cost is maintainer time. \\
\bottomrule
\end{tabularx}
\renewcommand{\arraystretch}{1.0}

% ══════════════════════════════════════════════════════════════════════════════
% FEATURE APPENDIX
% ══════════════════════════════════════════════════════════════════════════════

\newpage
\section*{Feature Appendix}

Defines the scope boundary for vox. Every feature is explicitly
categorized to prevent scope creep and force prioritization decisions upfront.

\subsection*{Shipped}

All features essential for the core value proposition are delivered and
available on PyPI.

\begin{enumerate}[nosep,leftmargin=2.5em]
  \featureitem{Stop hook notification}{When Claude Code finishes a task, synthesize and play a spoken summary of what happened. Shipped v0.2.0.}\label{feat:stop-hook}
  \featureitem{Permission-blocked notification}{When the agent waits for user approval, announce it. Uses the \texttt{Notification} hook with \texttt{permission\_prompt} subtype. Shipped v0.2.0.}\label{feat:permission}
  \featureitem{\texttt{/vox} command}{Enable or disable vox: \texttt{y} (chime notifications), \texttt{c} (continuous spoken summaries), \texttt{n} (off). Shipped v1.0.0.}\label{feat:notify}
  \featureitem{\texttt{/unmute} and \texttt{/mute} commands}{Toggle voice vs.\ chime. \texttt{/unmute} enables spoken notifications with optional voice selection. \texttt{/mute} switches to chime-only. Shipped v0.11.0.}\label{feat:speak}
  \featureitem{Chime audio}{Mood-aware audio tones for chime mode (\texttt{/mute}). Eight distinct signals pitch-shift to match session vibe (bright/neutral/dark). 24 total assets. Shipped v0.11.0.}\label{feat:chime}
  \featureitem{Provider auto-detection}{Detect best available TTS provider on first use. ElevenLabs > OpenAI > Polly > macOS \texttt{say} > Linux \texttt{espeak-ng}. Shipped v0.7.0.}\label{feat:autodetect}
  \featureitem{\texttt{/recap} command}{On-demand spoken summary of the last Claude response. Agent extracts key points and speaks them. Shipped v0.2.0.}\label{feat:recap}
  \featureitem{Plugin marketplace}{Listed in the Claude Code plugin marketplace for one-click install. Shipped v1.2.0.}\label{feat:marketplace}
  \featureitem{\texttt{voxd} audio daemon}{Dedicated audio server with WebSocket protocol. Separates synthesis/playback from MCP server for fast hook dispatch (observed on the order of a few milliseconds; not formally benchmarked). Playback queue, dedup, caching. Shipped v3.0.0.}\label{feat:voxd}
  \featureitem{Background music}{Vibe-driven instrumental music via ElevenLabs Music API. Adapts to session mood. \texttt{/music on|off}. The agent plays DJ --- a music-panel DJ-booth personality drives vibe-matched pools --- so vox is partly entertainment by design. Shipped v4.4.0.}\label{feat:music}
  \featureitem{Per-call API key passthrough}{Scope a single synthesis call to a specific provider API key for billing isolation. Four secure input paths. Shipped v4.3.0.}\label{feat:apikey}
  \featureitem{Config split}{Durable preferences (\texttt{vox.md}, tracked in git) separated from ephemeral session state (\texttt{vox.local.md}, gitignored). Shipped v4.7.5.}\label{feat:config-split}
  \featureitem{Continuous mode hooks}{UserPromptSubmit acknowledgment, SubagentStart/Stop announcements, SessionEnd farewell, PreCompact notification. Shipped v1.4.0.}\label{feat:continuous}
  \featureitem{Remote \texttt{voxd} connectivity}{Client-side env vars
  \texttt{VOXD\_HOST}, \texttt{VOXD\_PORT}, \texttt{VOXD\_TOKEN} override
  localhost discovery, enabling cross-host audio playback without SSH
  tunnels. Server-side \texttt{VOXD\_BIND} controls the bind address.
  Token auth secures the connection. Shipped v4.7.7.}\label{feat:remote-voxd}
  \featureitem{Three access surfaces}{The same TTS engine is reachable three
  co-equal ways: the \texttt{vox-server} MCP server (key \texttt{mic}) for
  agents, the \texttt{VoxClient}/\texttt{VoxClientSync} Python library, and the
  \texttt{vox} CLI. All three share one \texttt{voxd} daemon, so voice,
  caching, and playback are identical across surfaces. Unified over
  \texttt{voxd} since v3.0.0.}\label{feat:surfaces}
\end{enumerate}

\subsection*{Should Do / Next}

Deliberately deferred, on-roadmap scope --- not yet built, prioritized behind
adoption of the shipped feature set. Listed here so the deferral is a visible
positioning choice rather than an oversight.

\begin{enumerate}[nosep,leftmargin=2.5em]
  \featureitem{Broader downstream library adoption}{Grow the set of tools that use punt-vox as their TTS backend beyond langlearn-tts. The MCP server, Python library, and CLI surfaces are built; broader adoption is the deferred work, not new engineering.}\label{feat:downstream}
  \featureitem{Public recommendation and outreach}{Move from ecosystem/internal use to public recommendation --- marketplace promotion, documentation, and community outreach. Deferred until usage data validates the broadened engagement bet (see \faqref{faq:validation-plan}).}\label{feat:outreach}
  \featureitem{Codebase-aware audio programs}{Podcast and audiobook programs whose content is drawn from the codebase's own domain, the technology it uses, or fiction inspired by it --- so a developer can step back and decompress between sessions rather than burn out. This extends the DJ/entertainment side of vox beyond background music. Deliberately deferred roadmap, not shipped.}\label{feat:audio-programs}
\end{enumerate}

\subsection*{Won't Do}

Features explicitly excluded. Each exclusion is a deliberate positioning
choice.

\begin{enumerate}[nosep,leftmargin=2.5em]
  \featureitem{Voice input}{Listening to the developer and converting speech to commands. Different product, different architecture (ASR, wake word detection, continuous listening). vox speaks; it does not listen.}\label{feat:no-voice-input}
  \featureitem{Screen reader / accessibility narrator}{Continuous reading of terminal output for visually impaired users. Requires cursor tracking, ARIA-like semantics, and fundamentally different event architecture. This is valuable but is a different product.}\label{feat:no-screenreader}
  \featureitem{Custom voice training}{Training a TTS model on the user's voice or a custom voice. ElevenLabs supports this, but exposing it adds configuration complexity that contradicts the zero-config design principle.}\label{feat:no-custom-voice}
\end{enumerate}

% ══════════════════════════════════════════════════════════════════════════════
% REFERENCES
% ══════════════════════════════════════════════════════════════════════════════

\newpage
{\emergencystretch=1em
\printbibliography[title={References}]
}

\end{document}
